\chapter{\eh{electric-plug} Capa 1 --- Física (Physical Layer)}

\section{\eh{thinking-face} ¿Qué hace la capa física?}

\begin{concepto}
La \textbf{capa física} transmite \textbf{bits crudos} (ceros y unos) sobre un
medio físico: cable de cobre, fibra óptica o señales de radio. No sabe qué
significan los datos, solo se encarga de que lleguen correctamente.

\emoji{light-bulb} \textbf{Analogía:} Es como el sistema eléctrico de tu casa.
No sabe qué aparato vas a enchufar ni para qué lo usas. Solo garantiza que
la electricidad llega con la tensión correcta al tomacorriente.
\end{concepto}

\section{\eh{thread} Cable UTP --- Unshielded Twisted Pair}

\begin{concepto}
\textbf{UTP} (par trenzado no blindado) es el cable más común en redes LAN.
Tiene \textbf{4 pares de hilos de cobre}, cada par \textbf{trenzado} entre sí
dentro de una envoltura de plástico.

\emoji{light-bulb} \textbf{¿Por qué trenzado?} Al trenzar los cables se produce
cancelación de interferencias electromagnéticas (EMI) por inducción mutua.
Un par envía la señal en ambas polaridades; las interferencias externas afectan
igual a ambos hilos y al restar las señales, el ruido se cancela. Sin trenzado,
los cables actuarían como antenas captando ruido de motores, luces fluorescentes, etc.
\end{concepto}

\subsection{Categorías de cable UTP}

\begin{table}[h!]
\centering
\renewcommand{\arraystretch}{1.4}
\begin{tabular}{lllll}
\toprule
\textbf{Categoría} & \textbf{Velocidad máx.} & \textbf{Distancia} & \textbf{Frecuencia} & \textbf{Uso típico} \\
\midrule
Cat 5e  & 1 Gbps      & 100 m & 100 MHz  & Redes hogar/oficina básica \\
Cat 6   & 10 Gbps     & 55 m  & 250 MHz  & Redes corporativas \\
Cat 6a  & 10 Gbps     & 100 m & 500 MHz  & Data centers, acceso \\
Cat 7   & 10 Gbps     & 100 m & 600 MHz  & STP, blindado, industrial \\
Cat 8   & 40 Gbps     & 30 m  & 2000 MHz & Conexiones cortas en data center \\
\bottomrule
\end{tabular}
\caption{Categorías de cable UTP y sus especificaciones}
\end{table}

\begin{cuidado}
La categoría del cable determina su velocidad máxima. Si conectas Cat 5e
a un switch de 10 Gbps, el enlace negociará automáticamente a 1 Gbps.
\textbf{El eslabón más débil de la cadena limita a todos los demás.}
También importa la categoría de los jacks, patch panels y conectores: de nada
sirve Cat 6a si los jacks son Cat 5e.
\end{cuidado}

\subsection{Normas T568-A y T568-B}

\begin{concepto}
Las normas \textbf{ANSI/EIA/TIA T568-A} y \textbf{T568-B} definen el orden
exacto de los 8 hilos dentro del conector RJ-45. Son publicadas por TIA
(Telecommunications Industry Association).
\end{concepto}

\begin{center}
\begin{tikzpicture}[
    pin/.style={rectangle, draw=gray!60!white, minimum width=0.65cm, minimum height=1.4cm,
                font=\tiny\bfseries, align=center, text=textlight}
]
% T568-B
\node[font=\bfseries\small, text=textlight] at (2.6, 2.2) {T568-B (más común en México)};
\node[pin, fill=orange!60]          at (0.65,0.7) {};
\node[pin, fill=orange!90]          at (1.3, 0.7) {};
\node[pin, fill=green!50]           at (1.95,0.7) {};
\node[pin, fill=blue!50]            at (2.6, 0.7) {};
\node[pin, fill=blue!20]            at (3.25,0.7) {};
\node[pin, fill=green!80]           at (3.9, 0.7) {};
\node[pin, fill=brown!30]           at (4.55,0.7) {};
\node[pin, fill=brown!60]           at (5.2, 0.7) {};
\foreach \i/\n in {1/1,2/2,3/3,4/4,5/5,6/6,7/7,8/8}{
    \node[font=\footnotesize] at (\i*0.65-0.0, -0.1) {\n};
}
\node[font=\tiny, text=gray] at (0.65,0.0) {};
% Color labels
\node[font=\tiny, rotate=90, anchor=center] at (0.65,0.7)  {Bl-Nar};
\node[font=\tiny, rotate=90, anchor=center] at (1.3, 0.7)  {Nar};
\node[font=\tiny, rotate=90, anchor=center] at (1.95,0.7)  {Bl-Ver};
\node[font=\tiny, rotate=90, anchor=center] at (2.6, 0.7)  {Azul};
\node[font=\tiny, rotate=90, anchor=center] at (3.25,0.7)  {Bl-Az};
\node[font=\tiny, rotate=90, anchor=center] at (3.9, 0.7)  {Verde};
\node[font=\tiny, rotate=90, anchor=center] at (4.55,0.7)  {Bl-Caf};
\node[font=\tiny, rotate=90, anchor=center] at (5.2, 0.7)  {Café};
\end{tikzpicture}
\end{center}

\begin{table}[h!]
\centering
\renewcommand{\arraystretch}{1.3}
\begin{tabular}{clll}
\toprule
\textbf{Pin} & \textbf{T568-B} & \textbf{T568-A} & \textbf{Función (10/100)} \\
\midrule
1 & Blanco-Naranja & Blanco-Verde  & TX+ (transmisión positivo) \\
2 & Naranja        & Verde         & TX- (transmisión negativo) \\
3 & Blanco-Verde   & Blanco-Naranja & RX+ (recepción positivo) \\
4 & Azul           & Azul          & (no usado en 10/100) \\
5 & Blanco-Azul    & Blanco-Azul   & (no usado en 10/100) \\
6 & Verde          & Naranja       & RX- (recepción negativo) \\
7 & Blanco-Café    & Blanco-Café   & (no usado en 10/100) \\
8 & Café           & Café          & (no usado en 10/100) \\
\bottomrule
\end{tabular}
\caption{Orden de pines T568-A vs T568-B}
\end{table}

\begin{ejemplo}
\textbf{Cable straight-through (conexión directa):}
Ambos extremos usan el mismo estándar (T568-B). Conecta dispositivos de
diferente tipo: PC \textrightarrow{} Switch, Switch \textrightarrow{} Router.

\textbf{Cable crossover (cruzado):}
Un extremo T568-A, el otro T568-B. Los pares TX y RX se invierten.
Se usa para conectar dispositivos del mismo tipo: PC \textrightarrow{} PC,
Switch \textrightarrow{} Switch.

\emoji{light-bulb} \textbf{Hoy ya casi no se usa el crossover:} los switches
y NIC modernos tienen \textbf{Auto-MDI/MDIX} --- detectan automáticamente
si deben cruzar o no.

\textbf{Gigabit Ethernet (1000BASE-T)} usa los 4 pares (8 hilos) para
transmisión simultánea en ambas direcciones (full-duplex en cada par).
\end{ejemplo}

\section{\eh{optical-disk} Fibra Óptica}

\begin{concepto}
La fibra óptica transmite \textbf{pulsos de luz} en lugar de electricidad.
Ventajas: inmune a interferencia electromagnética, sin pérdida de señal
hasta kilómetros, y velocidades de Tbps.

\emoji{light-bulb} \textbf{Analogía:} Si el cable UTP es como hablar por
teléfono de lata con hilo (señal eléctrica que se degrada con la distancia),
la fibra es como usar un láser a través de un tubo de cristal perfecto:
la luz viaja casi sin pérdidas.
\end{concepto}

\begin{table}[h!]
\centering
\renewcommand{\arraystretch}{1.4}
\begin{tabular}{lll}
\toprule
\textbf{Característica} & \textbf{Multimodo (MMF)} & \textbf{Monomodo (SMF)} \\
\midrule
Diámetro del núcleo & 50 o 62.5 µm    & 8--10 µm \\
Fuente de luz       & LED / VCSEL      & Láser diodo \\
Distancia máxima    & hasta 550 m      & hasta 100+ km \\
Velocidad típica    & hasta 10 Gbps    & hasta 100 Gbps+ \\
Costo               & Menor            & Mayor \\
Color del conector  & Naranja o Aqua   & Amarillo \\
Uso                 & Dentro de edificios & Entre ciudades, ISP, Data Centers \\
\bottomrule
\end{tabular}
\caption{Fibra multimodo vs monomodo}
\end{table}

\begin{lstlisting}
   Multimodo                    Monomodo
   ┌─────────────────┐          ┌─────────────────┐
   │ núcleo 50/62µm  │ ←muchos  │ núcleo 8-10µm   │ ←un solo
   │ rayos de luz →  │  modos   │ rayo de luz →   │  modo
   └─────────────────┘          └─────────────────┘
   Dispersión modal → límite    Sin dispersión → largas distancias
   de distancia
\end{lstlisting}

\section{\eh{satellite-antenna} Wi-Fi --- Estándares IEEE 802.11}

\begin{concepto}
Wi-Fi usa \textbf{ondas de radio} para transmitir datos sin cables.
Cada generación de 802.11 mejora velocidad, distancia, o eficiencia.
\end{concepto}

\begin{table}[h!]
\centering
\renewcommand{\arraystretch}{1.4}
\begin{tabular}{lllll}
\toprule
\textbf{Estándar} & \textbf{Nombre} & \textbf{Frecuencia} & \textbf{Velocidad máx.} & \textbf{Año} \\
\midrule
802.11b  & Wi-Fi 1  & 2.4 GHz           & 11 Mbps   & 1999 \\
802.11g  & Wi-Fi 3  & 2.4 GHz           & 54 Mbps   & 2003 \\
802.11n  & Wi-Fi 4  & 2.4 / 5 GHz       & 600 Mbps  & 2009 \\
802.11ac & Wi-Fi 5  & 5 GHz             & 3.5 Gbps  & 2013 \\
802.11ax & Wi-Fi 6/6E & 2.4 / 5 / 6 GHz & 9.6 Gbps  & 2019 \\
\textbf{802.11be} & \textbf{Wi-Fi 7} & \textbf{2.4/5/6 GHz} & \textbf{46 Gbps} & \textbf{2024} \\
\bottomrule
\end{tabular}
\caption{Evolución de los estándares Wi-Fi}
\end{table}

\begin{moderno}
\textbf{\emoji{rocket} Wi-Fi 7 (802.11be, 2024)} introduce:
\begin{itemize}
    \item \textbf{Multi-Link Operation (MLO):} un dispositivo puede transmitir
          en múltiples bandas \textit{simultáneamente} (ej: 5 GHz y 6 GHz al mismo
          tiempo). Reduce latencia y duplica throughput efectivo.
    \item \textbf{Canales de 320 MHz:} el doble que Wi-Fi 6E (160 MHz).
    \item \textbf{4K-QAM:} modulación más densa, más bits por símbolo.
    \item \textbf{Latencia $<$1 ms:} clave para gaming, VR/AR, y control industrial.
    \item Teórico: 46 Gbps (en práctica con clientes reales: 5--15 Gbps).
\end{itemize}
\end{moderno}

\section{\eh{toolbox} Cableado Estructurado}

\begin{concepto}
Un sistema de \textbf{cableado estructurado} es la infraestructura de cables
y componentes instalada en un edificio de forma organizada según estándares
internacionales. Está diseñado para durar 15--25 años y soportar cualquier
tecnología de red futura.

La norma principal: \textbf{ANSI/EIA/TIA-568} (EE.UU.) y su equivalente
internacional \textbf{ISO/IEC 11801}.
\end{concepto}

Los \textbf{6 subsistemas} de un sistema de cableado estructurado:

\begin{enumerate}
    \item \textbf{Cableado horizontal:} del armario de telecomunicaciones
          al área de trabajo. Máximo 90 m de cable permanente.
    \item \textbf{Área de trabajo:} jacks RJ-45 en la pared, patch cords
          del jack al dispositivo (PC, teléfono IP, etc.).
    \item \textbf{Armario de telecomunicaciones (TR):} patch panels, switches,
          organizadores de cable. Uno por piso o zona.
    \item \textbf{Cableado backbone (vertical):} une los armarios entre sí
          y con la sala de equipos. Usa fibra óptica o UTP Cat 6a+.
    \item \textbf{Sala de equipos (ER):} servidores, routers, firewalls,
          equipos de core. Generalmente en sótano o piso dedicado.
    \item \textbf{Entrada de servicios:} donde el ISP o carrier llega
          al edificio. Demarc point (punto de demarcación).
\end{enumerate}

\begin{center}
\begin{tikzpicture}[
    box/.style={rectangle, draw=gray!50!white, rounded corners, minimum width=2.8cm,
                minimum height=0.8cm, font=\small, align=center, fill=blue!30!black, text=textlight},
    arr/.style={-{Stealth}, thick, draw=gray!60!white}
]
\node[box, fill=purple!35!black] (isp)   at (0, 4)   {ISP / Carrier};
\node[box, fill=orange!35!black] (er)    at (0, 2.8)  {Sala de Equipos\\(Subsistema 5)};
\node[box, fill=yellow!25!black] (back)  at (0, 1.6)  {Backbone Vertical\\(Subsistema 4)};
\node[box, fill=green!35!black]  (tr)    at (0, 0.4)  {Armario TR\\(Subsistema 3)};
\node[box, fill=cyan!35!black]   (horiz) at (0, -0.8) {Cable Horizontal\\(Subsistema 1)};
\node[box, fill=gray!40!black]   (wa)    at (0, -2.0) {Área de Trabajo\\(Subsistema 2)};
\node[box, fill=red!35!black]    (ent)   at (3.5,2.8) {Entrada de Servicios\\(Subsistema 6)};
\draw[arr] (isp.south) -- (er.north);
\draw[arr] (er.south)  -- (back.north);
\draw[arr] (back.south)-- (tr.north);
\draw[arr] (tr.south)  -- (horiz.north);
\draw[arr] (horiz.south)-- (wa.north);
\draw[arr] (ent.west)  -- (er.east);
\end{tikzpicture}
\end{center}

\begin{cuidado}
\textbf{Límites del cableado horizontal:}
\begin{itemize}
    \item Máx. 90 m de cable permanente (de patch panel a jack de pared)
    \item Máx. 10 m de patch cords (5 m en el TR + 5 m en el área de trabajo)
    \item Canal total máx: \textbf{100 m}
\end{itemize}
Pasarte de 100 m causa errores de señal, retransmisiones y velocidad degradada.
\end{cuidado}
